\subsection{Sprint 0}

Na Sprint 0, meu trabalho começou pela preparação dos repositórios e pela idealização da infraestrutura. Organizamos a base do código e avaliamos a criação de um \textit{mirror} dos repositórios como contingência para eventuais falhas no servidor da AGES. Em paralelo, rascunhei o uso de Docker para melhorar a \ac{dx} e iniciei sua integração no \textit{backend}, enquanto esboçávamos uma primeira proposta para os serviços e a arquitetura em nuvem necessários.

No eixo de pesquisa, levantamos materiais de referência para os AGES I e II nas tecnologias de \textit{backend} com NestJS e \textit{frontend} com Next.js, além de aprofundar o estudo dos serviços AWS que deveriam compor a solução. Também iniciei o aprendizado sobre testes no NestJS para sustentar práticas de qualidade desde o início. Houve ainda alinhamentos técnicos com os demais AGES III sobre ORMs, o que ajudou a consolidar diretrizes para a camada de persistência.

Além do que pensamos para a infraestrutura do projeto, o restante do time focou em entregar os protótipos das telas no \textit{Figma}. A única ressalva é que o foco nas telas foi tão grande que acabou adiando a modelagem do banco de dados para a próxima \textit{sprint}. Dado que o banco poderia vir a ser a parte mais complexa da nossa aplicação, algo que o escopo levantado com os clientes já indicava, acabamos sendo consideravelmente prejudicados.
\\

\fbox{%
    \parbox{0.95\linewidth}{%
        \setlength{\parindent}{15pt}
            \itshape O início deste projeto veio acompanhado de grande expectativa para mim, já que a AGES III era a etapa que eu mais desejava cursar. A experiência positiva no projeto anterior me motivou a entrar nessa nova fase com entusiasmo e com a intenção de ser um bom colega de equipe.

No entanto, logo nesta primeira \textit{sprint}, deparei-me com um cenário de fricções inesperadas entre os colegas, o que me surpreendeu bastante. Ao tentar sugerir o uso da arquitetura e de tecnologias de \textit{backend} que haviam funcionado bem em minha AGES anterior, fui confrontado sob a justificativa de que minhas escolhas não seriam adequadas. Além disso, quando procurei corrigir alguns trechos de código desses mesmos colegas, novos conflitos surgiram.

    Essas situações acabaram me desmotivando em certa medida. Ainda assim, por acreditar na importância de colaborar e entregar resultados ao projeto, decidi concentrar meus esforços na área de infraestrutura e evitar atritos desnecessários, priorizando a continuidade do trabalho em equipe.
    }%
}
